\textbf{Phần 1. Trắc nghiệm lựa chọn}
\setcounterpageref{ex}{0}

\begin{ex} %[1D8V29]
Tính kết quả của phép tính: $32 + 5$.
\choice
{$36$}
{$38$}
{\True $37$}
{$35$}
\loigiai{
Cộng hàng đơn vị: $2 + 5 = 7$. Giữ nguyên hàng chục là 3. Vậy $32 + 5 = 37$.
\textbf{Đáp số:} $37$.\par
\textbf{Đáp án của câu hỏi là C}
}
\end{ex}

\begin{ex} %[1D8V29]
Tính kết quả của phép tính: $64 + 3$.
\choice
{$66$}
{\True $67$}
{$68$}
{$65$}
\loigiai{
Cộng hàng đơn vị: $4 + 3 = 7$. Giữ nguyên 6 chục. Vậy $64 + 3 = 67$.
\textbf{Đáp số:} $67$.\par
\textbf{Đáp án của câu hỏi là B}
}
\end{ex}

\begin{ex} %[1D8V29]
Tính kết quả của phép tính: $81 + 7$.
\choice
{$87$}
{$89$}
{$86$}
{\True $88$}
\loigiai{
Cộng hàng đơn vị: $1 + 7 = 8$. Giữ nguyên 8 chục. Vậy $81 + 7 = 88$.
\textbf{Đáp số:} $88$.\par
\textbf{Đáp án của câu hỏi là D}
}
\end{ex}

\begin{ex} %[1D8V29]
Điền số thích hợp vào ô trống: $43 + \dots = 48$.
\choice
{\True $5$}
{$4$}
{$6$}
{$3$}
\loigiai{
Muốn tìm số hạng chưa biết, lấy: $48 - 43 = 5$.
\textbf{Đáp số:} $5$.\par
\textbf{Đáp án của câu hỏi là A}
}
\end{ex}

\begin{ex} %[1D8V29]
Điền số thích hợp vào ô trống: $\dots + 4 = 79$.
\choice
{$74$}
{$76$}
{\True $75$}
{$73$}
\loigiai{
Muốn tìm số hạng chưa biết, lấy: $79 - 4 = 75$.
\textbf{Đáp số:} $75$.\par
\textbf{Đáp án của câu hỏi là C}
}
\end{ex}

\begin{ex} %[1D8V29]
Tính giá trị của biểu thức: $52 + 3 + 4$.
\choice
{$58$}
{\True $59$}
{$60$}
{$57$}
\loigiai{
Thực hiện: $52 + 3 = 55$; $55 + 4 = 59$.
\textbf{Đáp số:} $59$.\par
\textbf{Đáp án của câu hỏi là B}
}
\end{ex}

\begin{ex} %[1D8V29]
Tính giá trị của biểu thức: $70 + 2 + 6$.
\choice
{$77$}
{$79$}
{$76$}
{\True $78$}
\loigiai{
Thực hiện: $70 + 2 = 72$; $72 + 6 = 78$.
\textbf{Đáp số:} $78$.\par
\textbf{Đáp án của câu hỏi là D}
}
\end{ex}

\begin{ex} %[1D8V29]
Bạn Lan có 24 cái kẹo, mẹ cho thêm 5 cái kẹo nữa. Hỏi Lan có tất cả bao nhiêu cái kẹo?
\choice
{\True $29$}
{$28$}
{$30$}
{$27$}
\loigiai{
Lan có tất cả: $24 + 5 = 29$ (cái kẹo).
\textbf{Đáp số:} $29$.\par
\textbf{Đáp án của câu hỏi là A}
}
\end{ex}

\begin{ex} %[1D8V29]
Lớp 1A có 31 học sinh, có thêm 4 học sinh mới chuyển đến. Hỏi lúc này lớp 1A có bao nhiêu học sinh?
\choice
{$34$}
{$36$}
{\True $35$}
{$33$}
\loigiai{
Số học sinh lớp 1A có lúc này là: $31 + 4 = 35$ (học sinh).
\textbf{Đáp số:} $35$.\par
\textbf{Đáp án của câu hỏi là C}
}
\end{ex}

\begin{ex} %[1D8V29]
Bác nông dân nuôi 42 con gà, bác mua thêm 6 con gà nữa. Hỏi đàn gà nhà bác có tất cả bao nhiêu con gà?
\choice
{$47$}
{\True $48$}
{$49$}
{$46$}
\loigiai{
Đàn gà có tất cả: $42 + 6 = 48$ (con gà).
\textbf{Đáp số:} $48$.\par
\textbf{Đáp án của câu hỏi là B}
}
\end{ex}

\begin{ex} %[1D8V29]
Tìm số tự nhiên $x$ thỏa mãn: $x - 3 = 65$.
\choice
{$67$}
{$69$}
{$66$}
{\True $68$}
\loigiai{
Ta có: $x = 65 + 3 = 68$.
\textbf{Đáp số:} $68$.\par
\textbf{Đáp án của câu hỏi là D}
}
\end{ex}

\begin{ex} %[1D8V29]
Tìm số tự nhiên $y$ thỏa mãn: $y - 6 = 91$.
\choice
{\True $97$}
{$96$}
{$98$}
{$95$}
\loigiai{
Ta có: $y = 91 + 6 = 97$.
\textbf{Đáp số:} $97$.\par
\textbf{Đáp án của câu hỏi là A}
}
\end{ex}

\textbf{Phần 2. Trắc nghiệm đúng sai}
\setcounterpageref{ex}{0}

\begin{ex} %[1D8V29]
Xét tính đúng hoặc sai của các khẳng định toán học sau đây:
\choiceTF
{\True Vườn nhà Mai có 53 cây hoa hồng, bố trồng thêm 5 cây hoa hồng nữa. trong vườn có tất cả $58$ cây hoa hồng.}
{Kết quả của bài toán «Một sợi dây dài 41 cm, nối thêm một đoạn dây dài 7 cm. Sợi dây mới dài bao nhiêu xăng-ti-mét» là $49\text{ cm}$.}
{Kết quả của bài toán «Tìm một số biết rằng lấy số đó trừ đi 4 thì được 52» là $57$.}
{\True Hai số có hiệu bằng 7, số bé là 82. số lớn là số $89$ .}
\loigiai{
\begin{itemize}
    \item Ý a) Đúng: Số cây hoa hồng là: $53 + 5 = 58$ (cây).
\textbf{Đáp số:} $58$.
    \item Ý b) Sai: Kết quả đúng là $48$, không phải $49$. Chiều dài sợi dây mới là: $41 + 7 = 48$ (cm).
\textbf{Đáp số:} $48$.
    \item Ý c) Sai: Kết quả đúng là $56$, không phải $57$. Số cần tìm là: $52 + 4 = 56$.
\textbf{Đáp số:} $56$.
    \item Ý d) Đúng: Số lớn là: $82 + 7 = 89$.
\textbf{Đáp số:} $89$.
\end{itemize}
\begin{itemize}
    \item a) \textbf{Đúng}.
    \item b) \textbf{Sai}.
    \item c) \textbf{Sai}.
    \item d) \textbf{Đúng}.
\end{itemize}
}
\end{ex}

\begin{ex} %[1D8V29]
Xét tính đúng hoặc sai của các bài toán và phép tính sau đây:
\choiceTF
{Kết quả của bài toán «Số 45 cộng với số nào để được 49» là $5$.}
{\True Kết quả của bài toán «Số 71 cộng với số nào để được 77» là $6$.}
{\True Kết quả của bài toán «Một cửa hàng buổi sáng bán được 63 chiếc bút, buổi chiều bán được thêm 5 chiếc bút. Cả ngày cửa hàng bán được bao nhiêu chiếc bút» là $68$.}
{Kết quả của bài toán «Tìm một số biết rằng lấy số đó trừ đi 4 thì được 52» là $57$.}
\loigiai{
\begin{itemize}
    \item Ý a) Sai: Kết quả đúng là $4$, không phải $5$. Số đó là: $49 - 45 = 4$.
\textbf{Đáp số:} $4$.
    \item Ý b) Đúng: Số đó là: $77 - 71 = 6$.
\textbf{Đáp số:} $6$.
    \item Ý c) Đúng: Cả ngày bán được: $63 + 5 = 68$ (chiếc bút).
\textbf{Đáp số:} $68$.
    \item Ý d) Sai: Kết quả đúng là $56$, không phải $57$. Số cần tìm là: $52 + 4 = 56$.
\textbf{Đáp số:} $56$.
\end{itemize}
\begin{itemize}
    \item a) \textbf{Sai}.
    \item b) \textbf{Đúng}.
    \item c) \textbf{Đúng}.
    \item d) \textbf{Sai}.
\end{itemize}
}
\end{ex}

\textbf{Phần 3. Trả lời ngắn}
\setcounterpageref{ex}{0}

\begin{ex} %[1D8V29]
Tính: $30 + 4 + 5$.
\shortans{39}
\loigiai{
\textbf{Đáp số:} $39$\par
Thực hiện: $30 + 4 = 34$; $34 + 5 = 39$.
\textbf{Đáp số:} $39$.
}
\end{ex}

\begin{ex} %[1D8V29]
Tìm số tự nhiên $a$ biết: $a + 2 = 86 + 1$.
\shortans{85}
\loigiai{
\textbf{Đáp số:} $85$\par
Vế phải: $86 + 1 = 87$. Ta có: $a + 2 = 87 \Rightarrow a = 87 - 2 = 85$.
\textbf{Đáp số:} $85$.
}
\end{ex}

\begin{ex} %[1D8V29]
So sánh: $51 + 6$ và $52 + 4$. Giá trị của biểu thức lớn hơn bằng bao nhiêu?
\shortans{57}
\loigiai{
\textbf{Đáp số:} $57$\par
$51 + 6 = 57$; $52 + 4 = 56$. Giá trị lớn hơn là $57$.
\textbf{Đáp số:} $57$.
}
\end{ex}

\begin{ex} %[1D8V29]
Hùng có 33 viên bi, Dũng cho Hùng thêm 6 viên bi. Sau đó Hùng làm mất 2 viên bi. Hỏi Hùng còn lại bao nhiêu viên bi?
\shortans{37}
\loigiai{
\textbf{Đáp số:} $37$\par
Sau khi nhận thêm Hùng có: $33 + 6 = 39$ (viên). Sau khi làm mất Hùng còn: $39 - 2 = 37$ (viên bi).
\textbf{Đáp số:} $37$.
}
\end{ex}

\begin{ex} %[1D8V29]
Băng giấy màu xanh dài 23 cm, cắt thêm 4 cm nữa để dài bằng băng giấy màu đỏ. Hỏi băng giấy màu đỏ dài bao nhiêu xăng-ti-mét?
\shortans{27}
\loigiai{
\textbf{Đáp số:} $27$\par
Băng giấy màu đỏ dài: $23 + 4 = 27$ (cm).
\textbf{Đáp số:} $27$.
}
\end{ex}

\begin{ex} %[1D8V29]
Tính: $92 + 1 + 2 + 3$.
\shortans{98}
\loigiai{
\textbf{Đáp số:} $98$\par
Thực hiện: $92 + 1 = 93$; $93 + 2 = 95$; $95 + 3 = 98$.
\textbf{Đáp số:} $98$.
}
\end{ex}
